\begin{table}[!htbp]\centering
\caption{Income's incremental contribution to cubic-score OLS fit}
\label{tab:premium-cubic-income-contribution}
\small
\begin{tabular}{lrrrr}\hline\hline
Outcome & $R^2$ without income & Full $R^2$ & Income $\Delta R^2$ & Income share (\%) \\ \hline
Premium institution & 0.2180 & 0.2493 & 0.0313 & 12.5 \\
Premium field & 0.0947 & 0.0978 & 0.0031 & 3.2 \\
Either (or both) & 0.1706 & 0.1729 & 0.0023 & 1.3 \\
Both & 0.1721 & 0.1783 & 0.0062 & 3.5 \\
Science & 0.0145 & 0.0164 & 0.0019 & 11.4 \\
Law & 0.0330 & 0.0347 & 0.0017 & 4.8 \\
Engineering / Manufacturing / Construction & 0.0426 & 0.0440 & 0.0014 & 3.2 \\
Medicine+ & 0.0410 & 0.0422 & 0.0012 & 2.8 \\
\hline\hline\end{tabular}
\begin{minipage}{0.98\textwidth}\footnotesize
\textit{Notes:} Income share is the increase in ordinary unadjusted $R^2$ from adding nine income-decile indicators, divided by full-model $R^2$. Both models retain all nine total-degree-three math/verbal terms and admission-year indicators, using the same 565,040 students. This is the unique incremental contribution after score/year controls; shared explanatory variation is not allocated to income. The 5,105 admission-2026 records remain. All fit measures are in-sample and descriptive.
\end{minipage}
\end{table}
